\chapter{Limits and known problems}
\label{ch:limits}

This chapter lists what hdl-vet does not do, and what is known to go wrong. Every
item was checked against the code of this version. Items with a transcript were run.
Items that say so were not checked by a test.

\section{What the sign-off does not see}

\begin{itemize}
\item \textbf{Logic.} The sign-off compares the structure (instances, ports,
  parameters, connections) and the diagnostics of the compiler. A raw edit that changes
  only logic passes without an \code{expect}. The first transcript below changes
  \code{count_o + 1'b1} into \code{count_o + 2'd2} in \code{timer}, and the sign-off
  says nothing.

  Use the operations where you can. They never write such an edit. If you
  use \apiname{Vet.splice}, you are responsible for the logic. The same holds for
  declarations, because \apiname{Module.add_wire} has no structural expectation.
\item \textbf{Partly undriven signals.} A signal can lose the driver of only some of
  its bits. This can happen through an operation that is not a checked operation (here
  \apiname{Instance.disconnect}), or through a splice. The sign-off passes such a run.

  The second transcript disconnects the interrupt of the second timer. Bit 1 of
  \code{irq_o} of \code{periph_wrap} has no driver, and the top module reads it. The
  sign-off accepts the run. Checked operations see this case (chapter~\ref{ch:operations},
  \file{remove\_timer.py}).
\item \textbf{Operations that cancel each other.} If one operation undoes another
  operation of the same script, the sign-off fails. An example is
  \apiname{Module.add_wire} and then \apiname{Module.remove_wire} of the same name. The
  check of the first operation does not find its result. The third transcript shows it.
  Remove both operations from the script.
\item \textbf{Gate primitives, user-defined primitives and checkers.} The sign-off does
  not compare them. A splice that changes \code{and g1 (z, a, b)} into
  \code{and g1 (z, a, a)} passes. hdl-vet writes the output tree, with exit status 0.
  The gate primitive was run. The other two follow from the same rule and were not checked
  by a test.
\item \textbf{More that the dangling check does not see.} The dangling check does not see
  these things:
  \begin{itemize}
  \item a parameter change that switches off a driver in a generate branch
  \item an inout port
  \item a new reader of a signal that had no driver already in the baseline
  \end{itemize}
  These cases were not checked by a test. The specification of the check lists them.
\item \textbf{One configuration.} The sign-off sees the design with the defines of the
  file list. An operation refuses a list that holds a conditional compilation directive,
  because it cannot check another configuration.
\end{itemize}

\lstinputlisting[style=script]{examples/scripts/logic_change.py}
\lstinputlisting[style=script]{examples/scripts/partial.py}
\lstinputlisting[style=script]{examples/scripts/cancel.py}

\lstinputlisting[style=console]{examples/console/limits.txt}

\section{What the operations refuse}

hdl-vet refuses to edit these things, and it prints a message. You can edit them by hand,
or with a raw splice.

\begin{itemize}
\item instances in generate blocks and instance arrays, and instances that come from a
  macro
\item hierarchical paths with an index
\item connection lists with positional or \code{.*} connections, and parameter lists
  with positional overrides
\item modules with a non-ANSI port list (for port operations), and module headers with
  a macro
\item ports and parameters inside conditional compilation directives
\item modules that are not instantiated under the top module
\item modules that come from an added file. hdl-vet never edits an added file.
\item interface ports. \apiname{Vet.punch_port} refuses them and
  \apiname{Instance.replace_module} refuses a module that has one. There is no operation
  that adds an interface port.
\item a replacement (\apiname{Instance.replace_module}) by a module whose port has another
  direction, width or type than the port that it replaces. There is no option to accept
  such a difference.
\item type parameters. \apiname{Instance.replace_module} does not carry an override of a
  type parameter to the new module. This holds even if the new module has a type parameter
  of that name. The message says that the new module has \emph{no overridable parameter}.
  Edit such an instance by hand.
\end{itemize}

\section{Known problems}

\begin{itemize}
\item \textbf{Escaped names.} \apiname{Module.remove_instance} cannot remove an instance
  with an escaped identifier, such as \texttt{\textbackslash u/esc}. If you give the name
  with the backslash, hdl-vet skips the call without a message, although the instance
  exists. If you give the name without the backslash, hdl-vet applies the call, but the
  sign-off fails with
  \code{unexpected change: instance-removed}.
\item \textbf{Names that the operations do not check.} \apiname{Module.add_port} and
  \apiname{Module.add_param} do not check the name that you give. They write
  \code{"a b"} or a keyword into the text, and the sign-off fails with errors of the
  compiler or with \code{the edited design cannot be loaded}. \apiname{Module.add_wire}
  does check its name.
\item \textbf{Time.} An operation that reads a list that an earlier operation edited
  loads the whole design again. On a large design, a script with many such operations takes
  a long time. The load count rose after a \code{connect} that follows an
  \code{add_port}. The times were not measured.
\item \textbf{Paths after a replacement.} After \apiname{Instance.replace_module}, a
  hierarchical operation (\apiname{Vet.punch_port}, \apiname{Vet.propagate_param}) still
  resolves its path against the original design. It edits the old module, and the sign-off
  reports the result as errors or failed checks. Do such operations before the
  replacement, or in another run.
\item \textbf{Scopes and shared modules.} hdl-vet refuses \apiname{Vet.punch_port} and
  \apiname{Vet.propagate_param} in a scope whose module the design instantiates more than
  once. The scope takes the name of a module, not an instance path. The check for
  references that reach outside a scope does not see a name in code that the compiler does
  not elaborate. Its message names only one reference when a block adds several.
\item \textbf{Added files.} A relative file name in \apiname{Vet.add_file} starts from the
  current folder, not from the script. Use \code{Path(__file__).resolve().parent} as in the
  example scripts.
\item \textbf{Saved sessions of an older version.} A session that you saved before
  hdl-vet recorded the nested file lists protects only the top-level file list when you
  resume it. With such a session, \opt{--save} onto a nested file list outside the input
  tree overwrites that list.
\item \textbf{Replacing a folder.} \opt{--force} with \opt{-o} replaces any folder that is
  not part of the input tree. This includes a folder with your own files in it. The rules
  are in chapter~\ref{ch:cli}.
\item \textbf{Old output trees.} An output tree that an older version wrote has no list
  of files in \file{hdl-vet.json}. \opt{--force} refuses it when the folder comes from the
  project file.
\end{itemize}
